\section{Rim, strain and calibrated mechanical feedback}
\label{sec:want-physical}
\textbf{Source binding.} This chapter lifts the geometry, strain, energy,
reporter and event models of WANTWOMBAN Operator I 1.1 (W1, physical PDF
pages 5--10, 13--17 and 20) into the R11 exact expression strategy. Source
equations are retained with their hypotheses; additional bounds below are
deductions for this subversion. No specimen, calibration series or photonic
gate performance is inferred from a mathematical expression.

The useful physical connection is explicit: a rim changes displacement;
displacement changes aperture geometry and gauge resistance; measurements
qualify the next completed event; the next command can change the physical
state. This supplies an instrument interface to the exact seed. It does not
yet implement the seed's arithmetic in an optical device.

\subsection{Geometry and a finite, declared mechanical realization}
Using W1 equations (4.1)--(4.7), choose a sheet radius $R>0$, positive
thicknesses $t,t_r$, rim height $h_r\ge0$, and a specified finite rib family:
\begin{align}
 \Omega_d&=\{(r,\vartheta,z):0\le r\le R,\ -t\le z\le0\},\\
 \Omega_r&=\{(r,\vartheta,z):R-t_r\le r\le R,\ 0\le z\le h_r\},\\
 \Omega_g&=(\Omega_d\cup\Omega_r\cup\bigcup_j\Omega_{{\rm rib},j})
                 \setminus\Omega_{\rm removed},\qquad
 m_g=\int_{\Omega_g}\rho_m\,dV .
\end{align}
Here $0<t_r\le R$, the angular coordinate covers a full turn, overlaps are
counted once, and removed material includes any actual slit or notch.
The word-decoded reference point $X$ has physical units only after a scale,
frame and origin are declared. Its deformed position is
$x(X,t)=X+u(X,t)$. A code for log-radius and a mechanical displacement are
different types even when both are carried by one-bit word streams.

For the small-gradient thermoelastic model,
\begin{equation}
 \epsilon=\tfrac12(\nabla u+\nabla u^T),\quad
 \epsilon_{\rm th}=\alpha_T(T-T_0)I,\quad
 \sigma=\mathbb C:(\epsilon-\epsilon_{\rm th}),\quad
 \rho_m\ddot u=\nabla\!\cdot\sigma+f .
\end{equation}
Displacement constraints, surface tractions, initial displacement and velocity,
material tensor and temperature field are part of this definition.
Positive elastic energy requires the appropriate symmetry and positivity of
$\mathbb C$ on the admissible strain space. A circular outline does not imply
isotropic material. A finite rotation with large displacement gradients is
outside this small-gradient constitutive model even if strain is small.

R11's finite-vector \code{ODE} is not a continuum PDE solver. To obtain a
closed finite profile, declare dimensionless admissible basis fields
$\phi_1,\ldots,\phi_m$ satisfying homogeneous displacement constraints and set
\begin{equation}
 u_m(X,t)=\sum_{i=1}^m q_i(t)\phi_i(X),\qquad
 B_i=\operatorname{sym}\nabla\phi_i .
\end{equation}
The coordinates $q_i$ have length units. For fixed geometry/material and a
declared damping matrix $D_m$, the weak-form coefficients are
\begin{align}
 M_{ij}&=\int_{\Omega_g}\rho_m\phi_i\!\cdot\!\phi_j\,dV,
 &K_{ij}&=\int_{\Omega_g}B_i:\mathbb C:B_j\,dV,\\
 (F_{\rm ext})_i&=\int_{\Omega_g}\phi_i\!\cdot f\,dV
                  +\int_{\Gamma_N}\phi_i\!\cdot\bar t\,dA,
 &(F_{\rm th})_i&=\int_{\Omega_g}B_i:\mathbb C:\epsilon_{\rm th}\,dV,\\
 M\ddot q+D_m\dot q+Kq&=F_{\rm ext}+F_{\rm th}.&&
\end{align}
With $M$ invertible this becomes the first-order exact declaration
\begin{equation}
 \dot q=v_q,\qquad
 \dot v_q=M^{-1}(F_{\rm ext}+F_{\rm th}-D_mv_q-Kq).
\end{equation}
For positive mass density and linearly independent admissible basis fields,
$z^TMz=\int\rho_m|\sum_i z_i\phi_i|^2dV>0$ for $z\ne0$, so the inverse
exists. Coefficient integrals require explicit charts, integrable functions
and Jacobians; nested finite \code{INTEGRAL} nodes can represent a selected
chart formula. They are not replaced by an unspecified measurement oracle.
Nonhomogeneous or moving supports require a lifting field and its extra
inertial/forcing terms. The basis can be analytic or modal: no raster grid is
mandated. Exact integration of this finite model still does not establish
that its basis truncation represents every deformation of a real specimen.

\subsection{What the rim benefit proves}
The same rectangular section in two orientations has
\begin{equation}
 I_{\rm upright}=t_rh_r^3/12,\qquad
 I_{\rm flat}=h_rt_r^3/12,\qquad
 I_{\rm upright}/I_{\rm flat}=(h_r/t_r)^2
\end{equation}
when $h_r>0$. This is a local section identity, not a complete assembly
stiffness ratio. W1's stronger assembly statement is conditional.
\begin{proposition}[Compliance under a positive stiffness addition]
If supported $K_0$ is symmetric positive definite, $K_r$ is symmetric
positive semidefinite, and the coordinate/load convention is unchanged, then
\begin{equation}
 F^T(K_0+K_r)^{-1}F\le F^TK_0^{-1}F.
\end{equation}
\end{proposition}
\begin{proof}
Let $A=K_0^{-1/2}K_rK_0^{-1/2}\succeq0$. The eigenvalues of
$(I+A)^{-1}$ lie in $(0,1]$. Congruence by $K_0^{-1/2}$ and evaluation on
$F$ gives the result.
\end{proof}
Mass redistribution, changed supports, prestress and buckling do not
automatically supply such a $K_r$. Nor does the inequality bound every local
strain component. In the scalar example $q=F/K$ and $\epsilon_g=a_qq$,
\begin{equation}
 V_o\simeq-\frac{V_{\rm ex}GF\,a_q}{4K}F.
\end{equation}
Increasing $K$ reduces displacement and this force-to-voltage sensitivity.
Thus aperture stability, sensing response and collecting area must be
specified together. A conditional design objective, retained from W1, is
\begin{equation}
 \min_{\mathcal G}\frac{m_g}{m_0}
   +w_c\frac{F^TK(\mathcal G)^{-1}F}{C_0},\quad
 P_{\rm det}\ge P_{\min},\quad
 |\partial V_o/\partial F|\ge S_{\min} .
\end{equation}
Its positive normalization scales, allowed geometries and load cases are
inputs. No optimum or optical performance is asserted. If the kernel switches
$K$, the storage change $\tfrac12q^T(K_{\rm new}-K_{\rm old})q$ enters the
work ledger; a passive fixed-$K$ proof does not erase this term.

\subsection{Exact resistance and bridge arithmetic}
For a uniform trace segment with positive electrical resistivity and lengths,
W1 equations (6.1)--(6.4) give
\begin{align}
 R_e&=\rho_eL/A, &\frac{R'_e}{R_e}
   &=\frac{\rho'_e}{\rho_e}\frac{\lambda_{\parallel}}
                                  {\lambda_{\perp1}\lambda_{\perp2}},\\
 \frac{\Delta R_e}{R_e}&\simeq\frac{\Delta\rho_e}{\rho_e}
   +\epsilon_{\parallel}-\epsilon_{\perp1}-\epsilon_{\perp2},
 &\epsilon_a&=a^T\epsilon a\quad(\|a\|=1),\\
 R_{\rm grid}&=\sum_iR_i,
 &\frac{\Delta R_{\rm grid}}{R_{\rm grid}}
    &=\sum_i\frac{R_i}{R_{\rm grid}}\frac{\Delta R_i}{R_i}.
\end{align}
The finite-stretch ratio assumes homogeneous stretches aligned with the
segment axes; the following strain expression is a first-order approximation.
The weights in the final identity are the reference resistances, so it is
an exact finite-change identity for the same series segments. Under isotropic
uniaxial stress, $\epsilon_{\perp j}=-\nu\epsilon_{\parallel}$ and the
small-signal gauge factor is
$GF\simeq1+2\nu+(\Delta\rho_e/\rho_e)/\epsilon_{\parallel}$ where that
ratio is defined. Adhesive strain transfer, transverse response and any
crystal-frame piezoresistivity require their own constitutive inputs.

Declare $R_1,R_2$ the left upper/lower arms, $R_3,R_4$ the right upper/lower
arms, all positive, with $V_o=V_L-V_R$ and $V_{\rm ex}>0$. Then
\begin{equation}
 v:=V_o/V_{\rm ex}=\frac{R_2}{R_1+R_2}-\frac{R_4}{R_3+R_4}.
 \label{eq:want-bridge}
\end{equation}
For $R_1=R_2=R_3=R$ and $R_4=R(1+x)$,
\begin{equation}
 x>-1,\qquad v=-\frac{x}{4+2x},\qquad
 x=-\frac{4v}{1+2v},\qquad -\tfrac12<v<\tfrac12 .
 \label{eq:want-inverse}
\end{equation}
Substitution proves inverse composition exactly. The denominator never
vanishes on the physical open domain, and $dv/dx=-1/(2+x)^2<0$, proving a
one-to-one map onto the stated interval. With rational resistance ratios,
these are rational operator graphs; no series expansion is needed.
The first-order output is $-V_{\rm ex}GF\epsilon/4$, with this lead polarity.
Replacing the exact nonzero output by $-x/4$ has signed relative error $x/2$.
R11 therefore preserves the rational transfer instead of accumulating that
known linearization error in the seed.

For four small changes, $v\simeq(-x_1+x_2+x_3-x_4)/4$.
An axial $R_4$ and bonded transverse $R_3$ with equal gauge factors give
$v\simeq-GF(1+\nu)\epsilon/4$ under the stated uniaxial model. The
transverse element contributes strain signal. Common multiplicative scaling
of each individual divider's two arms cancels exactly in
\eqref{eq:want-bridge}; cancellation of other thermal changes is a separate
matched-arm condition.

An exact source-decimal calibration profile can retain
\begin{equation}
 x=GF_0\epsilon+c_T\Delta T+c_{\epsilon T}\epsilon\Delta T
            +c_2\epsilon^2+c_0 .
\end{equation}
After exact bridge inversion, define $A=c_2$,
$B=GF_0+c_{\epsilon T}\Delta T$ and $C=c_T\Delta T+c_0-x$.
The strain is a selected root of $A\epsilon^2+B\epsilon+C=0$.
For $A=0,B\ne0$ it is $-C/B$; for $A\ne0$ the real candidates require
$B^2-4AC\ge0$ and are $(-B\pm\sqrt{B^2-4AC})/(2A)$.
A calibrated strain interval and root-selection rule are mandatory. If
$B+2A\epsilon$ has one strict sign throughout that interval, there is at
most one root there. Existence still requires the measured value to lie in
the model's range. Rational coefficients admit R11's algebraic root type;
more general real coefficients require its separate real-root declaration
and domain obligations. A convenient root cannot be silently selected.

\subsection{Measurement uncertainty survives an exact inverse}
For an isothermal profile with calibrated $GF=g>0$,
\begin{equation}
 \widehat\epsilon(v,g)=-\frac{4v}{g(1+2v)},\quad
 \frac{\partial\widehat\epsilon}{\partial v}
    =-\frac4{g(1+2v)^2},\quad
 \frac{\partial\widehat\epsilon}{\partial g}=-\widehat\epsilon/g .
\end{equation}
The inverse becomes increasingly sensitive as $v$ approaches $-1/2$.
This is conditioning of the physical inference, not a floating-point defect.
If $v\in[v_-,v_+]\subset(-1/2,1/2)$ with $g$ fixed, monotonicity gives
the exact enclosure
\begin{equation}
 \epsilon\in\left[
  -\frac{4v_+}{g(1+2v_+)},\ -\frac{4v_-}{g(1+2v_-)}\right].
 \label{eq:want-strain-enclosure}
\end{equation}
Interval calibration quantities can be propagated with denominator-sign
certificates and any model-residual bound. First-order uncertainty instead
uses $u_\epsilon^2\simeq J\Sigma J^T$ and retains covariance; it is not an
exact enclosure for arbitrary nonlinear uncertainty.

A digitizer code $m$ and its scale $\Delta$ can be retained exactly, while
the physical voltage lies in a declared uncertainty set around $m\Delta$.
An ideal rounding quantizer contributes at most $\Delta/2$ inside its valid
range; real offset, gain, drift and noise require additional characterization.
The exact code is evidence of a code, not an exact reconstruction of every
analogue quantity that could have produced it. W1's design matrix with rows
$[1,\epsilon,\Delta T,\epsilon\Delta T,\epsilon^2]$ must have rank five
to identify five independent coefficients in the noise-free linear fit.
Repeating one condition cannot repair deficient rank.

\section{Energy, reporter memory and completed optical events}
\label{sec:want-energy}
\subsection{Separate supplied energy from internal transfers}
W1 equations (14.1)--(14.3) introduce a useful closure:
\begin{equation}
 \Delta(U_{\rm th}+E_{\rm elastic}+K_{\rm mech}+E_{\rm chem})
 =E_{\rm abs}+E_{\rm electrical}+W_{\rm ext}
                   -Q_{\rm bath}-E_{\rm other,out}.
 \label{eq:want-ledger}
\end{equation}
The stores must form a nonduplicating constitutive partition. Light already
reflected or transmitted is excluded from absorbed energy; luminescence or
electricity exported after absorption belongs in the outgoing terms.
Chemical and mechanical heating transfer previously stored or supplied
energy, rather than adding independent copies of absorbed energy.

Electrical interrogation contributes real heat. The exact arm powers are
\begin{equation}
 P_i=\frac{V_{\rm ex}^2R_i}{(R_1+R_2)^2}\ (i=1,2),\qquad
 P_i=\frac{V_{\rm ex}^2R_i}{(R_3+R_4)^2}\ (i=3,4).
\end{equation}
At balance $P_i=V_{\rm ex}^2/(4R)$ and $P_{\rm total}=V_{\rm ex}^2/R$.
Only the characterized fractions coupled to each thermal node enter that
node. A bridge heater transient can affect the same reporter intended to
witness optical absorption. The electrical source and acquisition duty cycle
therefore belong in the physical model, regardless of exact digital arithmetic.

For $N$ thermal nodes use constant $C_i>0$, reciprocal
$G_{ij}=G_{ji}\ge0$, and $G_{i0}\ge0$:
\begin{equation}
 C_i\dot T_i=P_i-G_{i0}(T_i-T_a)
                    -\sum_jG_{ij}(T_i-T_j).
 \label{eq:want-network}
\end{equation}
$P_i$ includes separately accounted heater, bridge, absorbed-to-heat and
internal dissipation terms. Let $C=\operatorname{diag}(C_i)$,
$L_{ii}=G_{i0}+\sum_{j\ne i}G_{ij}$, $L_{ij}=-G_{ij}$ for $i\ne j$,
and $b=P+(G_{i0}T_a)_i$. Then $\dot T=-C^{-1}LT+C^{-1}b$ is a finite
R11 \code{ODE}. For constant coefficients its exact solution is denoted
\begin{equation}
 T(t+h)=e^{-C^{-1}Lh}T(t)
               +\int_0^h e^{-C^{-1}L(h-s)}C^{-1}b\,ds .
\end{equation}
Here the matrix exponential is derived notation: each column is the solution
of the corresponding linear finite ODE with a unit-vector initial value.
It is not a new opaque primitive or a claim of constant-cost evaluation.
This form permits a singular $L$; no unjustified $L^{-1}$ is introduced.

For one node, held $P,T_a,C$ and $G>0$, and duration $h\ge0$,
\begin{equation}
 T^+=T_a+P/G+(T-T_a-P/G)e^{-Gh/C}.
 \label{eq:want-thermal-step}
\end{equation}
For $G=0$ the correct branch is $T^+=T+hP/C$. The exponential step is exact
for the constant-input model. If absorption, heater power or material
coefficients change during the interval, their specified time dependence
must be included or the interval must be split. Exact evaluation of a held
approximation does not make the physical input constant.

\subsection{Bounded material-state memory}
For a reporter fraction $\zeta\in[0,1]$, W1 uses
\begin{equation}
 \dot\zeta=k_+(1-\zeta)-k_-\zeta,\qquad k_+,k_-\ge0 .
\end{equation}
A candidate photoactivation contribution is
$k_{\rm ph}=\int Y_\lambda\sigma_\lambda\Phi_\lambda\,d\lambda$,
with flux per area, time and wavelength, cross-section $\sigma_\lambda$,
and dimensionless conversion yield $Y_\lambda$; the product integrates to a
rate. Other activation/recovery processes are separate terms. The source
does not supply a material formulation or measured rates.

For held rates and $k=k_++k_->0$,
\begin{equation}
 \zeta^+=\zeta_{\rm eq}+(\zeta-\zeta_{\rm eq})e^{-kh},\qquad
 \zeta_{\rm eq}=k_+/k;qquad
 k=0:\ \zeta^+=\zeta .
 \label{eq:want-reporter}
\end{equation}
\begin{proposition}[Reporter interval preservation]
For $h\ge0$, nonnegative held rates, and $\zeta\in[0,1]$,
equation \eqref{eq:want-reporter} produces $\zeta^+\in[0,1]$.
\end{proposition}
\begin{proof}
For $k>0$, $r=e^{-kh}\in(0,1]$ and
$\zeta^+=r\zeta+(1-r)\zeta_{\rm eq}$ is a convex combination of two
members of $[0,1]$. The zero-rate branch is the identity.
\end{proof}
For time-varying integrable nonnegative rates the variation-of-constants
solution, with $K(t)=\int_{t_0}^t(k_++k_-)ds$, is
\begin{equation}
 \zeta(t)=e^{-K(t)}\zeta(t_0)
           +\int_{t_0}^t e^{-[K(t)-K(s)]}k_+(s)\,ds .
\end{equation}
It is nonnegative; applying the same formula to $1-\zeta$ with $k_+$ and
$k_-$ exchanged proves the upper bound. This extends interval preservation
without pretending the held-rate formula covers arbitrary inputs.

Reporter state can generate a fluorescence, resistance or thermal readout
only through a declared transduction law and noise model. Recovery and prior
dose are part of state. Two readouts of a shared absorption event can be
correlated; they do not provide two independent copies of its energy.

The retained ideal thermochromic update is
\begin{equation}
 a^+=\begin{cases}
 1,&T^+\le T_{\rm low},\\
 0,&T^+\ge T_{\rm high},\\
 a,&T_{\rm low}<T^+<T_{\rm high},
 \end{cases}\qquad T_{\rm low}<T_{\rm high},\quad a\in\B .
\end{equation}
This is a sampled hysteresis profile. A continuous relay with crossings
inside an interval requires an explicit event convention. Its state is also
distinct from a measured darkness bit with threshold noise. If an exact
expression comparison is unresolved, an expression evaluator must retain a
pending branch; a physical acquisition policy separately handles uncertain
or absent observations.

\subsection{One event bit from more than one observable}
For primary/secondary evidence $X,Y$, timing $C_t$, validity $V$, saturation
$S$ and fault $F$, W1 equation (12.1) is
\begin{equation}
 e=X\land Y\land C_t\land V\land\neg S\land\neg F.
\end{equation}
All six inputs are bits. Geometry-qualified mode uses strain to qualify
aperture conditions; this is not a second detection of photon absorption.
Pulse mode requires an absorber-linked physical response in a declared
energy range. Reporter mode additionally carries dose/recovery history.
Photon-resolved mode requires a characterized quantum instrument. These are
different event profiles, not a ladder of performance established by a flag.

The source's nominal lag gate is
\begin{equation}
 C_t=\mathbf1\{|(t_Y-t_X)-\tau_{\rm lag}|\le\Delta_c/2\},
 \qquad\Delta_c\ge0.
\end{equation}
For bounded timestamp/lag uncertainties, write the nominal difference as
$\widehat d$ and assume $|d-\widehat d|\le\delta_t$ with
$\delta_t\ge0$. The following stronger decisions are immediate:
\begin{align}
 |\widehat d|+\delta_t\le\Delta_c/2&\ \Longrightarrow\ 
                \text{all admitted timings pass},\\
 |\widehat d|-\delta_t>\Delta_c/2&\ \Longrightarrow\ 
                \text{all admitted timings fail}.
\end{align}
The remaining case is timing-uncertain. These follow from the triangle and
reverse triangle inequalities. Replacing it by a valid zero would change the
meaning of the observation; an event bit can still be withheld while keeping
the uncertainty status in the measurement record.

\subsection{Joint statistics with explicit conditioning}
Let $G=C_t\land V\land\neg S\land\neg F$ denote the quality/timing event.
For a null hypothesis $H_0$ and $P(G\mid H_0)>0$, the full false-acceptance
probability is
\begin{align}
 p_{\rm FA}&=P(G\mid H_0)P(X=1,Y=1\mid H_0,G),\\
 P(X=1,Y=1\mid H_0,G)
  &=\alpha_X\alpha_Y+\operatorname{Cov}(X,Y\mid H_0,G),
 \quad\alpha_X=P(X=1\mid H_0,G).
\end{align}
Define $\alpha_Y$ with the same conditioning. If $P(G\mid H_0)=0$ then
$p_{\rm FA}=0$ without invoking undefined conditional probabilities. Within
the conditioned population the sharp event bounds are
\begin{equation}
 \max(0,\alpha_X+\alpha_Y-1)
 \le P(X=1,Y=1\mid H_0,G)\le\min(\alpha_X,\alpha_Y).
\end{equation}
The lower bound follows from $P(X\cup Y)\le1$; the upper follows from set
inclusion. Products require independence under these same conditions, which
quality selection can itself alter. The same reasoning applies to detection
under $H_1$. If $Y$ is determined by $X$, then $I(H;Y\mid X)=0$; copied
software evidence is not an additional independent physical observation.

Under independent stationary Poisson background streams and a coincidence
window of total width $\Delta_c$, the expected pair rate is
$r_Xr_Y\Delta_c$. Counting each primary trigger at most once instead gives
$r_X(1-e^{-r_Y\Delta_c})$. These are specified stochastic models, not
measured rates for this assembly. Likewise, after $N>0$ independent identical
null trials with zero accepted events, a one-sided confidence $1-\alpha$
bound would be $p_{\rm FA,upper}=1-\alpha^{1/N}$ for $0<\alpha<1$.
Solving $(1-p)^N=\alpha$ derives the expression. No trials were run here,
and no confidence bound has been earned by this formalization.

\section{Deformation-aware exclusion and separate energy bounds}
\label{sec:want-pruning}
W1 equation (15.2) is a useful logical pruning certificate. For a fixed unit
slab normal $n$, centre coordinate $c$ and half-width $w/2$, suppose every
relevant physical descendant has
$x=P+(X-P)+u+e_x$ with
$|n\cdot(X-P)|\le R_N$, $|n\cdot u|\le d_N$ and
$|n\cdot e_x|\le\delta_N$ throughout the certified interval. Then
\begin{equation}
 |n\cdot P-c|>w/2+R_N+d_N+\delta_N
 \quad\Longrightarrow\quad |n\cdot x-c|>w/2
 \label{eq:want-slit-prune}
\end{equation}
by the reverse triangle inequality. Equality is retained. The descendants,
time interval and complete transformed child set must be covered; a tree
depth alone is not a geometric bound. The slab normal is the across-width
direction in the aperture plane, not automatically the optical propagation
axis. Exclusion from either enclosing width/height slab excludes a rectangular
aperture, provided all represented physical points use the same plane/frame.

\subsection{A moving-aperture certificate}
For a moving plane normal $n(t)$ and centre coordinate $c(t)$, assume
\begin{equation}
 \|n(t)-n_0\|\le\eta_n,\quad |c(t)-c_0|\le\eta_c,\quad
 \|x(t)\|\le B_x,\quad w(t)\le w_{\max}.
\end{equation}
If the previous projected $R_N,d_N,\delta_N$ bounds use $n_0$, then
\begin{equation}
 |n_0\cdot P-c_0|>
 w_{\max}/2+R_N+d_N+\delta_N+\eta_nB_x+\eta_c
 \label{eq:want-moving-prune}
\end{equation}
certifies exclusion for the whole interval. Indeed
$|n(t)\cdot x-c(t)|\ge|n_0\cdot x-c_0|-\eta_nB_x-\eta_c$.
All coordinates and $B_x$ share one declared origin; a large arbitrary
origin can make this bound unnecessarily loose. A local fixed origin with
the correspondingly transformed $c(t)$ improves its usefulness without
changing the proof. Exact packing can remove numerical encoding error from
a declared coordinate, but it cannot set physical deformation or mounting
uncertainty to zero.

For the source's rounded log-radius and phase codes, a relative radial
rounding envelope is $2^{\Delta\rho/2}-1$, and nearest 2048-bin phase has
angular envelope $\pi/2048$. These are inherited representation-profile
errors. An exact symbolic coordinate need not be rounded into those codes,
but using the existing W64 code does not restore the omitted information.
Convert angular/radial errors to physical length at the relevant scale before
adding them to a length bound.

\subsection{Omitted heating with a signed error budget}
Geometric exclusion does not prove absence of diffraction, scattering or
thermal coupling. For the constant passive network
\eqref{eq:want-network}, compare two runs with identical coefficients,
initial temperature and bath inputs, differing by heating $\delta P$.
If $\delta P_i(t)\ge0$, W1's thermal omission statement is
\begin{equation}
 0\le\delta T_i(t)\le E_{\rm omit}(t)/C_i,\qquad
 E_{\rm omit}(t)=\int_{t_0}^t\sum_j\delta P_j(s)\,ds.
\end{equation}
The following extension also handles a signed heat-model difference.
\begin{proposition}[Absolute heat-omission enclosure]
For the same passive linear network, zero initial temperature difference
and integrable signed $\delta P$, define
$E_{\rm abs}(t)=\int_{t_0}^t\sum_j|\delta P_j(s)|ds$. Then
\begin{equation}
 |\delta T_i(t)|\le E_{\rm abs}(t)/C_i.
\end{equation}
\end{proposition}
\begin{proof}
The matrix $A=-C^{-1}L$ has nonnegative off-diagonal entries, so its transition
matrix $e^{A(t-s)}$ is entrywise nonnegative. For example choose $\gamma$
so $A+\gamma I$ is entrywise nonnegative; the convergent exponential series
proves $e^{A\tau}=e^{-\gamma\tau}e^{(A+\gamma I)\tau}\ge0$ for $\tau\ge0$.
Thus the zero-initial solution $z$ driven by $|\delta P|$ satisfies
$|\delta T|\le z$ componentwise and $z\ge0$.
Summing $C_i\dot z_i$ cancels internal conductances and yields
$d(\sum_iC_iz_i)/dt=\sum_i|\delta P_i|-\sum_iG_{i0}z_i
\le\sum_i|\delta P_i|$. Integration gives
$C_iz_i\le\sum_jC_jz_j\le E_{\rm abs}$, proving the bound.
\end{proof}
The exponential series is used in the proof, not introduced as a finite
wire-format opcode. Feedback-dependent changes to power, conductance,
capacity, aperture or initial state violate this comparison's common-model
hypotheses and require a coupled bound. An optical omitted-field budget is
separate because coherent field interference is not bounded merely by a
geometrical occupancy bit.

\subsection{Integrated physical benefit and remaining obligations}
The new exact seed can retain the rational bridge law and calibrated roots,
exponential reporter/thermal states, deformation certificates, typed
uncertainty records and joint-event decision without rounding those internal
expressions back into persistent state. These are specific improvements to
the representation and auditability of a photonic sensing/control system.
Rim stiffness, optical transfer, reporter sensitivity and channel statistics
remain selected physical models with specimen-specific parameters.

The causal lift consumes a completed measurement at $t_n$, computes the next
command from that state, holds or otherwise declares that command over the
next interval, and records the resulting observation at $t_{n+1}$. Its event
can then affect the next branch. Delaying the feedback this way closes the
recurrence without requiring an observation to precede its own cause.
It is conditional on a uniquely defined plant update and declared noise or
input realization. Causality alone is neither stability nor a bound on
tracking error. An exact symbolic model can be useful before an optical
arithmetic device exists; demonstrating that device is a different task.
